\documentclass[10pt,a4paper]{article}
\usepackage[T1]{fontenc}
\usepackage[utf8]{inputenc}
\usepackage[margin=1.9cm]{geometry}
\usepackage{graphicx}
\usepackage{hyperref}
\setlength{\parindent}{0pt}
\setlength{\parskip}{3pt}
\raggedbottom

% Replace the names before submission.
\newcommand{\ReportAuthors}{[Name Surname; Name Surname]}

% Put the two images at these paths. Missing files leave visible placeholders,
% so this source can still be compiled before the images are available.
\newcommand{\reportfigure}[3]{%
  \begin{center}
    \IfFileExists{#1}{%
      \includegraphics[width=0.96\linewidth,height=0.30\textheight,keepaspectratio]{#1}%
    }{%
      \fbox{\parbox[c][0.27\textheight][c]{0.92\linewidth}{%
        \centering\small Add image file\\[5pt]\texttt{\detokenize{#1}}}}%
    }\\[5pt]
    \refstepcounter{figure}\label{#3}%
    \begin{minipage}{0.96\linewidth}
      \small\textbf{Figure \thefigure.} #2
    \end{minipage}
  \end{center}
}

\begin{document}
\begin{center}
  {\Large\bfseries FairJack: On-chain Blackjack with a Staked House Pool\par}
  \vspace{3pt}
  {\small INF0422 -- Blockchain, Distributed and Decentralised Systems\\
  \ReportAuthors}
\end{center}

\section*{1. Aim and system model}
FairJack is an educational decentralised blackjack prototype. A Solidity
contract holds the game's authoritative state and ERC-20 tokens; a React
interface connects players, liquidity providers and application-level
validators through a wallet. The test token FJT is freely mintable, so it is
used only to demonstrate the protocol, not to represent real value. The
project combines an on-chain state machine, shared financial accounting and
multi-party commit--reveal randomness.

Liquidity providers deposit FJT into a house pool and receive shares. Their
share price changes as the house wins or loses. A player stakes a bet
$b\in[1,100]$ FJT, and the pool reserves $2b$ until the game ends. Thus an
ongoing game cannot make its maximum payout unavailable through a withdrawal.
Three registered validators, distinct from the player, are assigned to each
game; each must have at least 100 FJT of collateral available per active
game. Selection starts at an index derived from a hash of the player, game ID
and contract address, then scans eligible validators. It is deterministic,
not a source of unpredictable randomness.

\section*{2. Game protocol}
The player and each of the three validators first submit
$H(s,\mathrm{gameId},\mathrm{sender},\mathrm{contract})$, where $s$ is a private
32-byte secret and $H$ is Keccak-256. Binding the commitment to the game,
account and contract prevents reuse in another context. All four commits
must arrive within 15 minutes. A further 15-minute reveal phase opens only
after the fourth commit. The contract checks every reveal against its
commitment and, after all four reveals, hashes the secrets in a fixed order
to form the final seed (Fig.~\ref{fig:protocol}).

Cards are drawn from a 52-card deck by hashing the seed with a changing nonce;
already drawn cards are rejected. Two cards go to the player and two to the
dealer. The player may \texttt{hit} or \texttt{stand}; each valid hit renews a
five-minute action deadline. A bust ends the game immediately. After a stand,
the contract draws for the dealer until its score reaches at least 17, then
settles by comparing scores. A win pays the player $2b$ in total, a push
returns $b$, and a loss transfers the bet's economic value to the pool.

\newpage
\section*{3. Fairness mechanisms and failure handling}
Commit--reveal prevents a participant from changing its secret after seeing
other commitments. The pool's shares and reserved-liquidity check allocate
ordinary wins and losses to the stakers while protecting active payouts.
Anyone may call \texttt{claimTimeout} after a deadline: a validator that
misses a commit or reveal loses 100 FJT collateral, credited to the pool;
an otherwise compliant player gets the bet back and the game is cancelled.
If the player misses its commit, reveal or action, the bet is forfeited.
Cancellation and settlement release the game's liquidity reservation and
validator capacity once. The contract uses OpenZeppelin's \texttt{SafeERC20}
for token transfers and \texttt{ReentrancyGuard} on methods that transfer
tokens or settle a game.

\section*{4. Evaluation and limitations}
The repository includes Hardhat tests for commitment validation, reveal
ordering, card uniqueness, scoring, payout accounting, deadlines, slashing
and a complete game. These tests specify intended behaviours; this report
does not claim an independent security audit or a deployed Sepolia instance.
The web interface exposes pool operations, validator management and the game
flow (Fig.~\ref{fig:interface}). The source contains a local deployment script and instructions in
\texttt{README.md}; the same script has a Sepolia configuration.

There are important fairness limits. A participant who reveals last can
choose to abort after seeing the other secrets; slashing makes this costly
but does not remove the option. Once the seed is published, the remaining
deck can be computed, so a player can choose hits with knowledge unavailable
in physical blackjack. The dealer's initial cards are also recorded in
public state. Validator selection is predictable and identities are not
Sybil-resistant. The FJT token and fixed 18-decimal bet constants further
limit the design to a controlled demonstration. These points prevent a claim
that the current game is secure for real-money use.

\section*{5. Reproduction and AI declaration}
The project archive should include the source and \texttt{README.md}, which
gives the local setup, deployment and browser walkthrough. In brief, install
the root and frontend dependencies, start the Hardhat node, run
\texttt{npm run deploy:local}, then \texttt{npm run frontend}; the README lists
the required test accounts and wallet network. AI tools supported the
testing work throughout the project, including security-oriented cases and
checks of code usability, and suggested refactorings where clearer
implementations were possible. OpenAI Codex was also used to analyse the
repository and draft this report.

{\footnotesize
\textbf{References.} [1] FairJack source: \texttt{contracts/FairJackPool.sol},
\texttt{contracts/MockERC20.sol}, \texttt{test/}, and \texttt{README.md}
(submitted version). [2] OpenZeppelin Contracts, ERC-20, SafeERC20 and
ReentrancyGuard (dependencies used by the source).
}

\newpage
\section*{Figures}
\reportfigure{figures/protocol-flow.png}{FairJack game lifecycle: commitments,
reveals, card dealing, player decisions and final settlement.}{fig:protocol}

\vfill

\reportfigure{figures/fairjack-interface.png}{FairJack web interface showing
game progress alongside pool or validator status.}{fig:interface}

\end{document}
